%% mathstc-signe-parabole — signe de f(x) = -2(x-1)(x-3) lu sur la parabole.
%% Le coefficient -2 est négatif : la parabole est tournée vers le bas, donc f(x) > 0
%% ENTRE les racines 1 et 3, et f(x) < 0 à l'extérieur. Aucune table de signes :
%% la figure doit rester une image mentale.
%% Échelles : 1 unité d'abscisse -> 1,5 cm ; 1 unité d'ordonnée -> 0,55 cm.
\documentclass[tikz,border=4pt]{standalone}
\usepackage{liaisons}
\begin{document}
\begin{tikzpicture}[x=1.5cm,y=0.55cm,
  axe/.style={-{Stealth[length=5pt]}, line width=0.6pt},
  courbe/.style={solideA, line width=1.3pt, line join=round},
  gris/.style={black!55, line width=0.6pt},
  grad/.style={font=\scriptsize, inner sep=2pt},
  note/.style={font=\scriptsize, text=black!60, inner sep=1.5pt}]

%% --- accolades « carrées » (ouverture vers le haut / vers le bas) ----------
%% #1 = x gauche, #2 = x droit, #3 = ordonnée, #4 = légende
\newcommand\accobas[4]{%
  \draw[gris] ($(#1,#3)+(0,-0.08cm)$) -- ($(#1,#3)+(0,-0.22cm)$)
           -- ($(#2,#3)+(0,-0.22cm)$) -- ($(#2,#3)+(0,-0.08cm)$);
  \pgfmathsetmacro\xm{0.5*(#1+#2)}
  \draw[gris] ($(\xm,#3)+(0,-0.22cm)$) -- ($(\xm,#3)+(0,-0.34cm)$);
  \node[note, anchor=north] at ($(\xm,#3)+(0,-0.34cm)$) {#4};}
\newcommand\accohaut[4]{%
  \draw[gris] ($(#1,#3)+(0,0.08cm)$) -- ($(#1,#3)+(0,0.22cm)$)
           -- ($(#2,#3)+(0,0.22cm)$) -- ($(#2,#3)+(0,0.08cm)$);
  \pgfmathsetmacro\xm{0.5*(#1+#2)}
  \draw[gris] ($(\xm,#3)+(0,0.22cm)$) -- ($(\xm,#3)+(0,0.34cm)$);
  \node[note, anchor=south] at ($(\xm,#3)+(0,0.34cm)$) {#4};}

%% ================= axes et graduations ====================================
\draw[axe] (-0.35,0) -- (4.35,0) node[anchor=south east, inner sep=3pt, font=\small] {$x$};
\draw[axe] (0,-6.4) -- (0,2.9) node[anchor=south west, inner sep=3pt, font=\small] {$y$};
\node[grad, anchor=north east] at (-0.04,-0.1) {$O$};
\foreach \xv in {1,2,3,4} { \draw[line width=0.5pt] (\xv,0.14) -- (\xv,-0.14); }
\node[grad, anchor=north] at (2,-0.2) {$2$};
\node[grad, anchor=north] at (4,-0.2) {$4$};
\foreach \yv in {-6,-4,-2,2} {
  \draw[line width=0.5pt] (0.05,\yv) -- (-0.05,\yv);
  \node[grad, anchor=east, xshift=-2pt] at (0,\yv) {$\yv$}; }

%% ================= la parabole ============================================
%% arcs situés SOUS l'axe : trait normal
\draw[courbe] plot[domain=0.1:1, samples=30, smooth] (\x,{-2*(\x-1)*(\x-3)});
\draw[courbe] plot[domain=3:3.9, samples=30, smooth] (\x,{-2*(\x-1)*(\x-3)});
%% arc situé AU-DESSUS de l'axe : surligné
\draw[courbe, line width=2.6pt]
     plot[domain=1:3, samples=40, smooth] (\x,{-2*(\x-1)*(\x-3)});


%% ================= les deux racines =======================================
\fill (1,0) circle (2pt);
\fill (3,0) circle (2pt);
\node[grad, anchor=north] at (1,-0.25) {$1$};
\node[grad, anchor=north] at (3,-0.25) {$3$};

%% ================= signes, en accolades ===================================
\accohaut{1}{3}{2.15}{$f(x)>0$}
\accobas{0.1}{1}{-5.7}{$f(x)<0$}
\accobas{3}{3.9}{-5.7}{$f(x)<0$}

%% ================= rappel sur le coefficient dominant =====================
\node[note, anchor=north east, align=right, text width=2.6cm] at (4.4,2.75)
     {le coefficient $-2$\\tourne la parabole\\vers le bas};

%% ================= l'expression, sous la figure ===========================
\node[font=\small, text=solideA, anchor=north] at (2,-7.3) {$f(x)=-2(x-1)(x-3)$};
\end{tikzpicture}
\end{document}
